% 埃尔德什·帕尔（Paul Erdős）（综述）
% license CCBYSA3
% type Wiki

本文根据 CC-BY-SA 协议转载翻译自维基百科\href{https://en.wikipedia.org/wiki/Paul_Erd\%C5\%91s}{相关文章}。

\begin{figure}[ht]
\centering
\includegraphics[width=6cm]{./figures/04eef634120f8175.png}
\caption{} \label{fig_ARPR_1}
\end{figure}
保罗·埃尔德什（英语：Paul Erdős，匈牙利语：Erdős Pál [ˈɛrdøːʃ ˈpaːl]；1913 年 3 月 26 日—1996 年 9 月 20 日）是一位匈牙利数学家。他是 20 世纪最具产出力的数学家之一，也是提出最多数学猜想的人之一。\(^\text{[2][3]}\)埃尔德什在离散数学、图论、数论、数学分析、逼近论、集合论以及概率论等领域追求并提出了大量问题。\(^\text{[4]}\)他的大部分研究集中在离散数学领域，解决了其中许多长期未解的难题。他大力倡导并作出重要贡献的拉姆齐理论，研究的是在何种条件下秩序必然出现。总体而言，埃尔德什的工作更倾向于解决已存在的开放问题，而非开辟或探索全新的数学领域。他一生共发表约 1500 篇数学论文，这一数字至今无人能及。\(^\text{[5]}\)

他以“社会化数学”的实践方式以及古怪的生活方式闻名；《时代》杂志称他为“怪人中的怪人”。\(^\text{[6]}\)他坚信数学是一项社会性活动，因此过着漂泊不定的生活，唯一目的就是与其他数学家共同撰写论文。即使在晚年，他仍将清醒的每个时刻都献给数学研究；1996 年，他在华沙参加数学会议时去世。\(^\text{[7]}\)

埃尔德什与合作者们的大量共同成果，促成了“埃尔德什数”的诞生——它表示某位数学家与埃尔德什之间的最短合著路径中的“步数”。
\subsection{生平}
保罗·埃尔德什于 1913 年 3 月 26 日出生在奥匈帝国的布达佩斯，\(^\text{[8]}\)是安娜（婚前姓威廉，Wilhelm）与拉约什·埃尔德什（原姓英兰德，Engländer）\(^\text{[9][10]}\)的独子。他的两个姐姐在他出生前几天因猩红热去世，年仅三岁和五岁。\(^\text{[11]}\)他的父母都是犹太人，任职于高中教授数学。埃尔德什从小便对数学表现出极大的兴趣。由于父亲在 1914 年至 1920 年间作为奥匈帝国的战俘被囚禁于西伯利亚，\(^\text{[10]}\)他幼年时期主要由一位德国女家庭教师照顾，\(^\text{[12]}\)母亲不得不长时间工作以维持生计。父亲在战俘营中自学了英语，但许多单词的发音不正确。后来父亲教儿子说英语时，保罗也继承了这种发音，并在其一生中都保持这种独特的口音。\(^\text{[13]}\)

他通过阅读父母放在家中的数学书籍自学阅读。五岁时，只要告诉他一个人的年龄，他就能在脑中计算出这个人一生到目前为止大约活了多少秒。\(^\text{[12]}\)由于两位姐姐早逝，他与母亲关系极为亲密，据说直到上大学前，他们一直同睡一张床。\(^\text{[14][15]}\)

16 岁时，父亲向他介绍了两个将成为他终生热爱的主题——无穷级数与集合论。在高中时期，埃尔德什热衷于解答《KöMaL》（《中学数学与物理杂志》）每月刊登的数学问题，并成为其中的活跃解题者之一。\(^\text{[16]}\)

埃尔德什在 17 岁时赢得全国考试后，进入布达佩斯大学学习。当时，匈牙利实行“入学限额制”（numerus clausus），对犹太学生的大学录取人数有严格限制。\(^\text{[13][17]}\)到他 20 岁时，他已找到伯特兰猜想（Bertrand’s postulate）的证明。\(^\text{[17]}\)1934 年，年仅 21 岁的他获得数学博士学位。\(^\text{[17]}\)他的博士导师是李波特·费耶尔（Lipót Fejér），这位著名数学家也是约翰·冯·诺伊曼、乔治·波利亚（George Pólya）以及保尔·图兰（Pál Turán）的导师。他随后在英国曼彻斯特大学获得博士后研究职位，因为当时匈牙利的犹太人正受到专制政权的压迫。在那里，他结识了 G.H. 哈代（G. H. Hardy）和斯坦尼斯瓦夫·乌拉姆（Stanisław Ulam）。\(^\text{[13]}\)

由于身为犹太人，匈牙利对他而言已不再安全，他于 1938 年离开祖国，移居美国。\(^\text{[17]}\)在第二次世界大战期间，埃尔德什的多位亲属在布达佩斯遇难，其中包括两位姑母、两位叔叔以及他的父亲，只有母亲幸存下来。当时，埃尔德什正在美国普林斯顿高等研究院（Institute for Advanced Study）任职。\(^\text{[17][18]}\)然而，由于他不符合该机构的行为标准（他们认为他“粗野且不合常规”\(^\text{[13]}\)），他的奖学金并未像预期那样延长一年，而只续期了六个月。

传记作者保罗·霍夫曼（Paul Hoffman）形容他“或许是世界上最古怪的数学家”。\(^\text{[19]}\)埃尔德什成年后的大部分人生都过着提着手提箱四处漂泊的生活。除了 1950 年代几年的特殊时期（他因被指“同情共产主义者”而被禁止进入美国），他几乎一生都在从一个会议或研讨会赶往另一个会议。\(^\text{[19]}\)每到一个地方，他都期望主人为他提供住宿、饮食、洗衣以及一切生活所需，并帮他安排前往下一站的行程。\(^\text{[19]}\)

1943 年，乌拉姆离开威斯康星大学麦迪逊分校，前往新墨西哥州洛斯阿拉莫斯参与“曼哈顿计划”，与其他数学家和物理学家共事。他曾邀请埃尔德什加入该项目，但当埃尔德什表达了战后想回匈牙利的意愿后，邀请被撤回。\(^\text{[13]}\)

1996 年 9 月 20 日，83 岁的埃尔德什在华沙参加学术会议时突发心脏病去世。\(^\text{[20]}\)他的离世几乎与他自己理想的方式一致。他曾说：

“我希望能在讲课时去世，当我在黑板上完成一个重要证明时，听众中有人喊道：‘那一般情形呢？’我会转身微笑着说：‘我把那留给下一代吧。’然后我就倒下。”\(^\text{[20]}\)

埃尔德什终身未婚，也没有子女。\(^\text{[9]}\)他安葬于布达佩斯的犹太人科兹马街公墓，与父母长眠于同一处。\(^\text{[21]}\)他曾为自己的墓志铭建议写上：“我终于不再变笨了。”（匈牙利语：“Végre nem butulok tovább”）。\(^\text{[22]}\)

埃尔德什的名字中包含匈牙利字母“ő”（双上尖音符的 o），但常被误写为 Erdos 或 Erdös，原因“要么是错误，要么是出于印刷上的需要”。\(^\text{[23]}\)
\subsection{职业生涯}
1934 年，埃尔德什前往英国曼彻斯特任客座讲师。1938 年，他接受了美国新泽西州普林斯顿高等研究院（Institute for Advanced Study）的奖学金职位，这是他在美国的第一个学术任职，并持续了接下来的十年。在此期间，他与马克·卡茨（Mark Kac）和奥雷尔·温特纳（Aurel Wintner）合作发表了关于概率数论的重要论文，与保尔·图兰（Pál Turán）合作研究逼近论，并与维托尔德·胡列维奇（Witold Hurewicz）合作研究维数理论。\(^\text{[24]}\)尽管成果卓著，他的奖学金并未续期，埃尔德什不得不以“流浪学者”的身份在宾夕法尼亚大学、圣母大学、普渡大学、斯坦福大学和雪城大学等多所高校间辗转任职。他从未在一个地方长期停留，而是在世界各地的数学机构之间不断旅行，直至去世。

由于当时的“红色恐慌”（Red Scare）与麦卡锡主义运动，\(^\text{[25][26][27]}\)1954 年，美国移民与归化局拒绝为埃尔德什（匈牙利公民）签发重入美国的签证。\(^\text{[28]}\)当时他正在圣母大学任教，本可选择留在美国，但他决定收拾行李离开，虽然此后仍定期向美国移民局提出重新审查的申请。此后，他一度移居以色列，在耶路撒冷希伯来大学获得为期三个月的任职，随后又在以色列理工学院（Technion）获得“永久访问教授”的职位。
